\documentclass[12 pt, letterpaper]{hmcpset}
\usepackage{math-19-preamble}

\name{Jayden Thadani}
\class{MATH 142 --- Differential Geometry}
\assignment{Homework assignment \#2}
\duedate{2nd February 2024}

\begin{document}

\begin{problem}[Problem 1.]
	We say a function $\alpha : I \to \mathbb R^3$ is a \textit{regular curve} if, upon writing $\alpha(t) = (\alpha_1(t), \alpha_2(t), \alpha_3(t))$, there is no $t \in I$ such that
	\[
		\frac{d \alpha_1}{dt} \Big |_{t_0} = \frac{d \alpha_2}{dt} \Big |_{t_0} = \frac{d \alpha_3}{dt} \Big |_{t_0} = 0
	\]
	For a fixed $t_0 \in I$, the \textit{tangent line} to a regular curve is the straight line $\ell : \mathbb R \to \mathbb R^3$ given by
	\[
		\ell(u) = \left ( \alpha_1(t_0) + u \frac{d \alpha_1}{dt} \Big |_{t_0}, \alpha_2(t_0) + u \frac{d \alpha_2}{dt} \Big |_{t_0}, \alpha_3(t_0) + u \frac{d \alpha_3}{dt} \Big |_{t_0} \right ).
	\]
	Find the tangent line to the helix given by $\alpha(t) = (2 \cos t, 2 \sin t, t)$ for
	\begin{enumerate}
		\item $t_0 = 0$
		\item $t_0 = \pi/4$.
	\end{enumerate}
\end{problem}

\begin{solution}
	We can write out component functions $\alpha_1, \alpha_2, \alpha_3$ as
	\[
		\begin{gathered}
			\alpha_1(t) = 2 \cos t \\
			\alpha_2(t) = 2 \sin t \\
			\alpha_3(t) = t.
		\end{gathered}
	\]
	We can also write out their derivatives
	\[
		\begin{gathered}
			\frac{d \alpha_1}{dt} \Big |_{t_0} = - 2 \sin t_0 \\
			\frac{d \alpha_2}{dt} \Big |_{t_0} = 2 \cos t_0 \\
			\frac{d \alpha_3}{dt} \Big |_{t_0} = 1.
		\end{gathered}
	\] \\
	From these we can get the tangent lines for each $t_0$ just by calculating
	\begin{enumerate}
		\item
			\[
				\boxed{
					\ell(u) = (2, 2 u, u)
				}
			\] \\
			We get this by calculating
			\[
				\begin{split}
					\ell(u) & = \left ( \alpha_1(0) + u \frac{d \alpha_1}{dt} \Big |_0, \alpha_2(0) + u \frac{d \alpha_2}{dt} \Big |_0, \alpha_3(0) + u \frac{d \alpha_3}{dt} \Big |_0 \right ) \\
						& = ( 2 \cos 0 + u (-2 \sin 0), 2 \sin 0 + u (2 \cos 0), 0 + u) \\
						& = (2, 2u, u).
				\end{split}
			\]
		\item
			\[
				\boxed{
					\ell(u) = (\sqrt 2 (1 - u), \sqrt 2 (1 + u), \pi/4 + u).
				}
			\] \\
			We get this by calculating
			\[
				\begin{split}
					\ell(u) & = \left ( \alpha_1(0) + u \frac{d \alpha_1}{dt} \Big |_0, \alpha_2(0) + u \frac{d \alpha_2}{dt} \Big |_0, \alpha_3(0) + u \frac{d \alpha_3}{dt} \Big |_0 \right ) \\
						& = ( 2 \cos (\pi/4) + u (-2 \sin (\pi/4)), 2 \sin (\pi/4) + u (2 \cos (\pi/4)), \pi/4 + u) \\
						& = (\sqrt 2 - u \sqrt 2, \sqrt 2 + u \sqrt 2, \pi/4 + u).
				\end{split}
			\]
	\end{enumerate}
\end{solution}

\pagebreak

\begin{problem}[Problem 2]
	Let $\mathbf v = (1, 2, -3)$ and $\mathbf p = (0, -2, 1)$. Evaluate the following $1$-forms on the tangent vector $[\mathbf v, \mathbf p]$.
	\begin{enumerate}
		\item $y^2 \, dx$
		\item $z \, dy - y \, dz$
		\item $(z^2 - 1) \, dx - dy + x^2 \, dz$
	\end{enumerate}
\end{problem}

\begin{solution}
	For all of these, we can just calculate explicitly
	\begin{enumerate}
		\item 
			\[
				\boxed{
					y^2 \, dx ([\mathbf v, \mathbf p]) = 4.
				}
			\]\\
			We get this by
			\[
				\begin{split}
					y^2 \, dx([\mathbf v, \mathbf p]) & = (-2)^2 \, v_1 \\
									  & = 4 \cdot 1 \\
									  & = 4.
				\end{split}
			\]
		\item
			\[
				\boxed{
					(z \, dy - y \, dz)([\mathbf v, \mathbf p]) = -4.
				}
			\]\\
			We get this by
			\[
				\begin{split}
					(z \, dy - y \, dz)([\mathbf v, \mathbf p]) & = (1 \cdot v_2 - (-2) \, v_3) \\
										    & = (1 \cdot 2 - (-2)(-3)) \\
										    & = -4.
				\end{split}
			\]
		\item
			\[
				\boxed{
					((z^2 - 1) \, dx - dy + x^2 \, dz)([\mathbf v, \mathbf p]) = -2.
				}
			\]\\
			We get this by
			\[
				\begin{split}
					((z^2 - 1) \, dx - dy + x^2 \, dz)([\mathbf v, \mathbf p]) & = (1^2 - 1) \, v_1 - v_2 + 0^2 v_3 \\
												   & = 0 - 2 + 0 \\
												   & = -2.
				\end{split}
			\]
	\end{enumerate}
\end{solution}

\pagebreak

\begin{problem}[Problem 3]
	A $1$-form $\varphi : \mathrm T \mathbb R^3 \to \mathbb R$ is \textit{zero} at $\mathbf p$ if $\varphi([\mathbf v, \mathbf p]) = 0$ for all tangent vectors $[\mathbf v, \mathbf p] \in \mathrm T_{\mathbf p} \mathbb R^3$. Assume that $\varphi = df$ for a smooth function $f : \mathbb R^3 \to \mathbb R$. A point $\mathbf p \in \mathbb R^3$ at which the differential $\varphi = df$ is zero is called a \textit{critical point} for the function $f$.
	\begin{enumerate}
		\item Prove that $\mathbf p$ is a critical point of $f$ if and only if 
			\[
				\frac{\partial f}{\partial x} \Big |_{\mathbf p} = \frac{\partial f}{\partial y} \Big |_{\mathbf p} = \frac{\partial f}{\partial z} \Big |_{\mathbf p} = 0.
			\]
		\item Find all critical points of $f(x, y, z) = (1 - x^3) y + (1 - y^2) z$
	\end{enumerate}
\end{problem}

\begin{solution}
	\begin{enumerate}
		\item 
			First we will show that if $\mathbf p$ is a critical point, we have the desired equation. \\\\
			Note that $df([\mathbf v, \mathbf p])$ is just $\nabla f(\mathbf p) \cdot \mathbf v$. Thus, if $df([\mathbf v, \mathbf p])$ is zero for all tangent vectors in $\mathrm T_{\mathbf p} \mathbb R^3$, we must have $\nabla f(\mathbf p) \cdot \mathbf v = 0$ for all vectors $\mathbf v \in \mathbb R^3$. \\\\
			The only such vector is $\mathbf 0$, and thus, we must have $\nabla f(\mathbf p) = \mathbf 0$, or component-wise,
			\[
				\frac{\partial f}{\partial x} \Big |_{\mathbf p} = \frac{\partial f}{\partial y} \Big |_{\mathbf p} = \frac{\partial f}{\partial z} \Big |_{\mathbf p} = 0.
			\] \\
			Now we will show that if the equation holds, $\mathbf p$ is a critical point. Note that if the equation holds, we have
			\[
				\begin{split}
					\nabla f(\mathbf p) & = \left ( \frac{\partial f}{\partial x} \Big |_{\mathbf p}, \frac{\partial f}{\partial y} \Big |_{\mathbf p}, \frac{\partial f}{\partial z} \Big |_{\mathbf p} \right ) \\
							    & = (0, 0, 0) \\
							    & = \mathbf 0.
				\end{split}
			\] \\
			This gives us
			\[
				\begin{split}
					df([\mathbf v, \mathbf p]) & = \nabla f(\mathbf p) \cdot \mathbf v \\
								   & = \mathbf 0 \cdot \mathbf v \\
								   & = 0.
				\end{split}
			\]
		\item
			\[
				\boxed{
					\begin{gathered}
						\mathbf p_1 = \left ( 0, 1, \frac{1}{2} \right ) \\
						\mathbf p_2 = \left ( 0, -1, -\frac{1}{2} \right ).
					\end{gathered}
				}
			\] \\
			First, we compute $\nabla f$ to see what its zeroes are. For $\mathbf p = (x, y, z)$
			\[
				\begin{split}
					\nabla f(\mathbf p) & = \left ( \frac{\partial f}{\partial x} \Big |_{\mathbf p}, \frac{\partial f}{\partial y} \Big |_{\mathbf p}, \frac{\partial f}{\partial z} \Big |_{\mathbf p} \right ) \\
							    & = (3 x^2 y, 1 - x^3 - 2y z, 1 - y^2)
				\end{split}
			\] \\
			Since the critical points of $f$ are exactly the points where its gradient is zero, we can just solve for $\nabla f(\mathbf p) = \mathbf 0$.
			\[
				\begin{gathered}
					3 x^2 y = 0 \\
					1 - x^3 - 2y z = 0 \\
					1 - y^2 = 0.
				\end{gathered}
			\]
			Note that in order to solve this system $y = \pm 1$. Since $y \neq 0$, to solve the first equation we must have $x = 0$. This leaves us with two ways to solve the second equation:
			\[
				1 - 2z = 0 \text{ for } y = 1 \quad \text{ or } \quad 1 + 2z = 0 \text{ for } y = - 1.
			\] \\
			Thus, we have critical points $\mathbf p_1 = (0, 1, 1/2)$ and $\mathbf p_2 = (0, -1, -1/2)$.
	\end{enumerate}
\end{solution}

\pagebreak

\begin{problem}[Problem 4]
	For any three $1$-forms $\phi_i = \sum_j f_{ij} \, dx_j$, for $1 \le i \le 3$, prove
	\[
		\phi_1 \wedge \phi_2 \wedge \phi_3 = \begin{vmatrix}
			f_{11} & f_{12} & f_{13} \\
			f_{21} & f_{22} & f_{23} \\
			f_{31} & f_{32} & f_{33}
		\end{vmatrix} \, dx_1 \wedge dx_2 \wedge dx_3
	\]
\end{problem}

\begin{solution}
	We know already that
	\[
		\begin{split}
		\phi_2 \wedge \phi_3 & = (f_{21} \, f_{32} - f_{31} \, f_{22}) \, dx_1 \wedge dx_2 + (f_{22} \, f_{33} - f_{32} \, f_{23}) \, dx_2 \wedge dx_3 \\
				     & + (f_{23} \, f_{31} - f_{33} \, f_{21}) \, dx_3 \wedge dx_1 \\
		& = \sum_{j = 1}^3 (f_{2j} \, f_{3k} - f_{3j} f_{2k}) \, dx_j \wedge dx_k
		\end{split}
	\]
	where $k$ is $\sigma(j)$ where $\sigma$ is the cycle $(1 \, 2 \, 3)$. Now we only need to wedge an additional $\phi_3$ to this. \\\\
	Again, first we use linearity to get
	\[
		\begin{split}
			\phi_1 \wedge \phi_2 \wedge \phi_3 & = \left ( \sum_{i = 1}^3 f_{1 i} \, dx_i \right ) \wedge \left ( \sum_{j = 1}^3 (f_{2j} \, f_{3k} - f_{3j} f_{2k}) \, dx_j \wedge dx_k \right ) \\
							   & = \sum_{i = 1}^3 \sum_{j = 1}^3 f_{1 i} (f_{2j} \, f_{3k} - f_{3j} f_{2k}) \, dx_i \wedge dx_j \wedge dx_k.
		\end{split}
	\]
	By applying the alternating property, we know that any wedges with any two indices equal are $0$ and can cancel them out. This means we have to cancel out any terms where $i = j$ or $i = k$ ($j \neq k$). If $i = 1$, then $j$ can only be $2$, and so on --- for any choice of $i$, $j$ must be $\sigma(i)$. \\\\
	Writing out the three terms corresponding to $i = 1$, $i = 2$, and $i = 3$, we get
	\[
		\begin{split}
			\phi_1 \wedge \phi_2 \wedge \phi_3 & = f_{1 1} (f_{2 2} \, f_{3 3} - f_{3 2} f_{2 3}) \, dx_1 \wedge dx_2 \wedge dx_3 \\ 
							   & + f_{1 2} (f_{2 3} \, f_{3 1} - f_{3 2} f_{2 1}) \, dx_2 \wedge dx_3 \wedge dx_1 \\
							   & + f_{1 3} (f_{2 1} \, f_{3 2} - f_{3 1} f_{2 2}) \, dx_3 \wedge dx_1 \wedge dx_2.
		\end{split}
	\]
	Since all of these wedges are an even number of transpositions away from $dx_1 \wedge dx_2 \wedge dx_3$ they are all just $dx_1 \wedge dx_2 \wedge dx_3$ (the first one is equal, shift $dx_1$ to the left twice for the second, shift $dx_3$ to the right twice for the third). This allows us to write $\phi_1 \wedge \phi_2 \wedge \phi_3$ as
	\[
		(f_{1 1} (f_{2 2} \, f_{3 3} - f_{3 2} f_{2 3}) + f_{1 2} (f_{2 3} \, f_{3 1} - f_{3 2} f_{2 1}) + f_{1 3} (f_{2 1} \, f_{3 2} - f_{3 1} f_{2 2})) \, dx_1 \wedge dx_2 \wedge dx_3.
	\]
	which is just expanding along the first row of the determinant
	\[
		\begin{vmatrix}
			f_{11} & f_{12} & f_{13} \\
			f_{21} & f_{22} & f_{23} \\
			f_{31} & f_{32} & f_{33}
		\end{vmatrix} \, dx_1 \wedge dx_2 \wedge dx_3.
	\]

\end{solution}

\pagebreak

\begin{problem}[Problem 5]
	Let $f$ and $g$ be real-valued functions on $\mathbb R^2$. Prove that
	\[
		df \wedge dg = \begin{vmatrix}
			\dfrac{\partial f}{\partial x} & \dfrac{\partial f}{\partial y} \\
			\dfrac{\partial g}{\partial x} & \dfrac{\partial g}{\partial y}
		\end{vmatrix} \, dx \wedge dy.
	\]
\end{problem}

\begin{solution}
	Note that we can explicitly compute $df$ and $dg$ in terms of the partials of $f$ and $g$. That is, 
	\[
		df = \dfrac{\partial f}{\partial x} \, dx + \dfrac{\partial f}{\partial y} \, dy
	\]
	and
	\[
		dg = \dfrac{\partial g}{\partial x} \, dx + \dfrac{\partial g}{\partial y} \, dy.
	\]\\
	Now, to calculate their wedge product, we first use linearity to get
	\[
		\begin{split}
			df \wedge dg & = \left ( \dfrac{\partial f}{\partial x} \, dx + \dfrac{\partial f}{\partial y} \, dy \right ) \wedge \left ( \dfrac{\partial g}{\partial x} \, dx + \dfrac{\partial g}{\partial y} \, dy \right ) \\
				     & = \dfrac{\partial f}{\partial x} \, \dfrac{\partial g}{\partial x} \, dx \wedge dx + \dfrac{\partial f}{\partial y} \, \dfrac{\partial g}{\partial y} \, dy \wedge dy + \dfrac{\partial f}{\partial x} \, \dfrac{\partial g}{\partial y} \, dx \wedge dy + \dfrac{\partial f}{\partial y} \, \dfrac{\partial g}{\partial x} \, dy \wedge dx.
		\end{split}
	\]
	From here, we can apply the alternating property that $dx \wedge dx = dy \wedge dy = 0$ and $dx \wedge dy = - dy \wedge dx$ to get
	\[
		\begin{split}
			df \wedge dg & = \left ( \dfrac{\partial f}{\partial x} \, \dfrac{\partial g}{\partial y} - \dfrac{\partial f}{\partial y} \, \dfrac{\partial g}{\partial x} \right ) \, dx \wedge dy \\
				     & = \begin{vmatrix}
			\dfrac{\partial f}{\partial x} & \dfrac{\partial f}{\partial y} \\
			\dfrac{\partial g}{\partial x} & \dfrac{\partial g}{\partial y}
		\end{vmatrix} \, dx \wedge dy
		\end{split}
	\]
	as desired.
\end{solution}

\end{document}
